% problems.tex — \input dari report.tex
\chapter{Permasalahan Koordinasi}
\label{ch:problems}

Karena kedua lapis pemulihan beroperasi pada granularitas berbeda dan tidak berbicara
secara baku, interaksinya memunculkan lima \emph{titik gesekan}. Setiap masalah disajikan
dengan: konteks skenario, diagram \emph{sekuens} (pertukaran pesan antar-komponen),
diagram \emph{aktivitas} (alur keputusan), pseudokode, lalu pembahasan yang menelusuri
ketiganya. Warna merah menandai jalur kegagalan, hijau menandai jalur yang diinginkan.
Tiap masalah dikonkretkan dengan aplikasi \emph{ilustratif} yang sengaja
direka---\emph{PayGate}, \emph{SensorHub}, \emph{ArenaServer}, \emph{CoEdit}, dan
\emph{TopScore}---bukan produk nyata, melainkan contoh hipotetis untuk membumikan
persoalan.

\section{Ketidakcocokan Granularitas \emph{Liveness}}
\label{sec:p1}
Layanan \emph{webhook} pembayaran \emph{PayGate} menilai tiap transaksi melalui API
penilai \emph{fraud} eksternal lewat \emph{pool} Finch---pustaka klien HTTP Elixir yang
mengelola sekumpulan koneksi keluar yang dipakai ulang. Ketika API tersebut menggantung
(TCP setengah terbuka, tanpa \emph{timeout}), seluruh pekerja \emph{pool} terblokir dan
semua permintaan macet, tetapi \emph{listener} Cowboy---server HTTP Phoenix yang menerima
koneksi pada \emph{port}---tetap meng-\emph{accept} \emph{socket} sehingga \emph{probe} TCP
lulus dan Kubernetes mengira \emph{pod} sehat. Sebaliknya, \emph{probe} yang terlalu agresif
justru membunuh VM yang sebenarnya akan dipulihkan OTP.

\begin{figure}[H]\centering\includegraphics[width=\linewidth,height=0.9\textheight,keepaspectratio]{01-liveness-mismatch-sequence}
  \caption{Masalah 1 (sekuens): jalur sebuah permintaan ke \emph{PayGate}. Cowboy tetap
  menerima koneksi TCP dan pekerja Finch terblokir pada API \emph{fraud} yang menggantung;
  karena \emph{probe} hanya menguji penerimaan \emph{socket}, ia lulus meski 100\%
  pembayaran macet.}\label{fig:p1-seq}\end{figure}
\begin{figure}[H]\centering\includegraphics[width=\linewidth,height=0.9\textheight,keepaspectratio]{01-liveness-mismatch-activity}
  \caption{Masalah 1 (aktivitas): dua mode kegagalan rancangan \emph{probe}. \emph{Probe}
  dangkal melewatkan subsistem yang macet (negatif palsu); \emph{probe} agresif dengan
  \texttt{failureThreshold} rendah membunuh VM yang sebenarnya dapat dipulihkan OTP
  (positif palsu).}\label{fig:p1-act}\end{figure}
\begin{algorithm}[H]\caption{Masalah 1 --- \emph{probe} TCP naif.}\label{alg:p1}
\begin{algorithmic}[1]
\Function{NaiveProbe}{}
  \If{\Call{TcpAccept}{\textit{port}}} \State \Return \textit{HEALTHY} \Comment{meski \emph{pool} macet} \EndIf
\EndFunction
\end{algorithmic}\end{algorithm}

\noindent\textbf{Pembahasan.} Diagram sekuens (Gambar~\ref{fig:p1-seq}) menelusuri satu
transaksi: \emph{Client}~$\rightarrow$~Cowboy~$\rightarrow$~Finch~$\rightarrow$~API
\emph{fraud}, lalu API menggantung sehingga pekerja tak pernah dibebaskan. Pesan
\emph{probe} dari kubelet yang datang sesudahnya tetap dibalas ``\emph{ACK}'' oleh Cowboy
karena penerimaan \emph{socket} berada di jalur yang berbeda dari \emph{pool} yang macet.
Diagram aktivitas (Gambar~\ref{fig:p1-act}) memisahkan persoalan menjadi dua cabang
keputusan rancangan: cabang ``\emph{probe} dangkal'' berhenti pada ``\emph{port} terbuka?''
sehingga melewatkan kemacetan internal, sedangkan cabang ``\emph{probe} agresif'' bercabang
pada \texttt{failureThreshold} yang terlalu kecil dan mematikan VM saat gangguan sesaat.
Algoritma~\ref{alg:p1} memformalkan akarnya: vonis \textit{HEALTHY} pada baris~2 hanya
bergantung pada \textsc{TcpAccept}, tanpa sama sekali menyentuh keadaan \emph{pool} pekerja
atau \emph{supervision tree}---inilah celah semantik yang ditutup M1 dan M9.

\subsection{Ilustrasi Sederhana}
\label{subsec:p1-ilustrasi}
Aplikasi \emph{PayGate} memanggil API \emph{fraud} lewat \emph{connection pool} (sekumpulan
koneksi siap-pakai ke layanan luar). Langkahnya:
\begin{enumerate}[nosep,leftmargin=2em]
  \item API \emph{fraud} menggantung---tidak membalas dan tanpa batas waktu---sehingga satu
        per satu pekerja di \emph{pool} ikut menunggu selamanya, hingga semua pekerja habis
        dan setiap pembayaran macet.
  \item Pada saat yang sama, kubelet (agen Kubernetes di tiap node yang menjalankan
        pemeriksaan kesehatan) mengirim \emph{probe} TCP: ia hanya mengecek apakah \emph{port}
        4000 mau menerima koneksi.
  \item Karena server web (Cowboy) tetap menerima koneksi di \emph{port} itu---menerima
        koneksi adalah pekerjaan yang terpisah dari memprosesnya---\emph{probe} lulus, dan
        Kubernetes menyimpulkan \emph{pod} sehat.
\end{enumerate}
\noindent Akar masalahnya: \emph{probe} hanya bertanya ``\emph{port}-nya hidup?'', bukan
``aplikasinya masih mengerjakan sesuatu?''. Inilah yang diperbaiki M1 dan M9.

\section{Pemulihan Ganda}
\label{sec:p2}
Pada \emph{ingester} telemetri IoT \emph{SensorHub}, sebuah \emph{frame} cacat membuat
\emph{decoder}---sebuah \texttt{GenServer} (proses \emph{stateful} Elixir)---terus
\emph{crash}. Bila \texttt{max\_restarts} (jatah \emph{restart} yang dipantau supervisor)
terlalu longgar,
\emph{process} ber-\emph{crash-loop} tanpa terlihat oleh Kubernetes; bila diatur wajar,
habisnya jatah \emph{restart} mengeskalasikan kegagalan ke Kubernetes (\emph{CrashLoopBackOff})
secara sengaja.

\begin{figure}[H]\centering\includegraphics[width=\linewidth,height=0.9\textheight,keepaspectratio]{02-double-recovery-sequence}
  \caption{Masalah 2 (sekuens): perangkat mengirim ulang \emph{frame} cacat sehingga
  \emph{decoder} \emph{crash} berulang. Bila jatah \emph{restart} longgar, \emph{loop}
  tak terlihat oleh Kubernetes; bila wajar, supervisor menghentikan diri dan VM keluar
  sehingga Kubernetes mengambil alih.}\label{fig:p2-seq}\end{figure}
\begin{figure}[H]\centering\includegraphics[width=\linewidth,height=0.9\textheight,keepaspectratio]{02-double-recovery-activity}
  \caption{Masalah 2 (aktivitas): batas eskalasi kepemilikan. Selama dalam jatah, OTP
  yang memulihkan; saat jatah habis, kepemilikan berpindah ke Kubernetes.}\label{fig:p2-act}\end{figure}
\begin{algorithm}[H]\caption{Masalah 2 --- pemulihan tak terkoordinasi.}\label{alg:p2}
\begin{algorithmic}[1]
\Procedure{OnCrash}{\textit{child}}
  \State \Call{OTP.Restart}{\textit{child}} \Comment{max\_restarts longgar $\Rightarrow$ loop tak terlihat}
  \State \textbf{// VM tetap hidup; Kubernetes tidak diberi tahu}
\EndProcedure
\end{algorithmic}\end{algorithm}

\noindent\textbf{Pembahasan.} Diagram sekuens (Gambar~\ref{fig:p2-seq}) memperlihatkan
\emph{loop} berulang \emph{Device}~$\rightarrow$~\emph{Decoder}~$\rightarrow$~Supervisor:
perangkat mengirim ulang \emph{frame} yang sama, \emph{decoder} \emph{crash}, supervisor
menyalakannya ulang, lalu siklus berputar. Pada cabang ``jatah longgar'', tidak ada satu
pun pesan yang mengalir ke Kubernetes---itulah sebabnya kegagalan menjadi tak terlihat.
Diagram aktivitas (Gambar~\ref{fig:p2-act}) menempatkan keputusan kunci pada ``apakah
intensitas \emph{restart} melampaui jatah?''; selama \emph{tidak}, OTP yang memiliki
kegagalan, dan begitu \emph{ya}, kepemilikan berpindah ke Kubernetes melalui keluarnya VM.
Algoritma~\ref{alg:p2} sengaja ditulis ``naif'': baris~2 selalu me-\emph{restart} tanpa
memeriksa jatah, dan komentar baris~3 menyorot bahwa lapis orkestrator tak pernah
diberi tahu---persis ketiadaan kebijakan kepemilikan yang ditangani M2.

\subsection{Ilustrasi Sederhana}
\label{subsec:p2-ilustrasi}
Pada \emph{SensorHub}:
\begin{enumerate}[nosep,leftmargin=2em]
  \item Sebuah \emph{frame} (paket data) yang rusak dari perangkat membuat \emph{decoder}
        (\emph{process} yang menerjemahkan data masuk) \emph{crash}.
  \item Supervisor (\emph{process} pengawas yang otomatis menyalakan ulang anaknya saat
        \emph{crash}) langsung menghidupkannya kembali; perangkat mengirim ulang \emph{frame}
        rusak yang sama, \emph{decoder} \emph{crash} lagi, dihidupkan lagi---berputar terus.
  \item Selama VM BEAM tetap hidup, Kubernetes tak melihat masalah apa pun (\emph{probe}
        lulus) sehingga diam saja. Bila jatah \emph{restart} (\texttt{max\_restarts}) diatur
        wajar, supervisor akhirnya menyerah dan VM keluar, barulah Kubernetes me-\emph{restart}
        \emph{pod}.
\end{enumerate}
\noindent Akar masalahnya: tak ada kesepakatan lapis mana yang ``memiliki'' kegagalan ini,
sehingga bisa terjadi pemulihan berulang yang sia-sia. Ditangani M2 dan M8.

\section{\emph{Readiness} vs.\ Pembentukan \emph{Cluster}}
\label{sec:p3}
Pada \emph{backend} permainan \emph{ArenaServer}, tiap pertandingan adalah \emph{singleton}
global yang dikelola Horde (\emph{registry}/\emph{supervisor} terdistribusi dengan
\emph{handoff}). Saat \emph{rolling deploy}, sebuah \emph{pod} baru ditandai siap (hanya
karena \emph{port} terbuka) sebelum libcluster (pembentuk \emph{cluster} BEAM) menemukan
\emph{peer} dan Horde tersinkron;
lalu lintas tiba terlalu dini dan \emph{registry} masih kosong.

\begin{figure}[H]\centering\includegraphics[width=\linewidth,height=0.9\textheight,keepaspectratio]{03-readiness-cluster-sequence}
  \caption{Masalah 3 (sekuens): \emph{pod} menandai diri siap segera setelah \emph{port}
  terbuka, lalu Kubernetes merutekan permintaan pemain sebelum Horde tersinkron sehingga
  pencarian \emph{singleton} menemui \emph{registry} kosong.}\label{fig:p3-seq}\end{figure}
\begin{figure}[H]\centering\includegraphics[width=\linewidth,height=0.9\textheight,keepaspectratio]{03-readiness-cluster-activity}
  \caption{Masalah 3 (aktivitas): perbandingan kebijakan \emph{gating readiness}---``\emph{port}
  terbuka'' (terlalu dini) vs.\ ``\emph{cluster} tergabung dan Horde konvergen''
  (benar).}\label{fig:p3-act}\end{figure}
\begin{algorithm}[H]\caption{Masalah 3 --- \emph{readiness} terlalu dini.}\label{alg:p3}
\begin{algorithmic}[1]
\Procedure{OnPodStart}{}
  \State \Call{BindPort}{4000}; \Call{MarkReady}{} \Comment{libcluster/Horde belum konvergen}
\EndProcedure
\Procedure{OnRequest}{\textit{match}}
  \State \Call{Lookup}{\textit{match}} \Comment{$\rightarrow$ notFound / duplikat}
\EndProcedure
\end{algorithmic}\end{algorithm}

\noindent\textbf{Pembahasan.} Diagram sekuens (Gambar~\ref{fig:p3-seq}) menyusun urutan
waktu yang kritis: \emph{pod} mengikat \emph{port}, kubelet mengirim \emph{probe} TCP yang
langsung lulus, Kubernetes menambahkan \emph{pod} ke \emph{Endpoints}, lalu permintaan
pemain tiba---semua \emph{sebelum} pesan penemuan libcluster dan sinkronisasi Horde
selesai. Diagram aktivitas (Gambar~\ref{fig:p3-act}) mengontraskan dua kebijakan
\emph{gating}: jalur ``\emph{port} terbuka'' langsung menandai siap, sedangkan jalur yang
benar menunggu syarat gabungan ``\emph{cluster} tergabung \emph{dan} Horde konvergen''.
Algoritma~\ref{alg:p3} memisahkan dua prosedur: \textsc{OnPodStart} (baris~2) menandai
siap segera setelah \textsc{BindPort}, sementara \textsc{OnRequest} (baris~5) melakukan
\textsc{Lookup} yang mengembalikan \emph{notFound} atau bahkan mendaftarkan \emph{singleton}
duplikat---akibat langsung dari penandaan siap yang prematur, yang diperbaiki M4.

\subsection{Ilustrasi Sederhana}
\label{subsec:p3-ilustrasi}
\emph{ArenaServer} adalah \emph{backend} sebuah permainan daring \emph{real-time} berbasis
pertandingan---bayangkan game pertarungan arena (sejenis MOBA atau \emph{battle royale})
tempat banyak pemain bermain bersama dalam satu pertandingan. Tiap pertandingan yang sedang
berlangsung diwakili oleh satu \emph{process} di BEAM yang menyimpan keadaan permainan (skor,
posisi pemain, dan sebagainya).

Karena pemain berdatangan melalui \emph{load balancer}, dua pemain dalam pertandingan yang
sama bisa saja tersambung ke \emph{pod} yang berbeda. Agar keduanya tetap dilayani oleh
\emph{process} pertandingan yang sama, \emph{pod}-\emph{pod} itu harus saling mengenal dan
membentuk satu \emph{cluster} BEAM---\textbf{inilah mengapa \emph{peer} diperlukan}: tanpa
mengetahui rekan-rekannya, sebuah \emph{pod} tak bisa menemukan pertandingan yang sebenarnya
berjalan di \emph{pod} lain. \textbf{Mengapa perlu sinkronisasi}: daftar ``pertandingan~$X$
berjalan di \emph{pod}~$Y$'' harus diketahui seragam oleh semua \emph{pod}, dan Horde menjaga
daftar (\emph{registry}) ini tetap tersinkron sehingga tiap \emph{pod} dapat mengarahkan
pemain ke tempat yang benar. \textbf{Mengapa perlu \emph{rolling deploy}}: memperbarui game
(versi baru, perbaikan bug) dilakukan bertahap agar layanan tetap hidup tanpa
\emph{downtime}; \emph{pod} lama diganti \emph{pod} baru sedikit demi sedikit---dan di sinilah
\emph{pod} baru bermunculan, yang harus bergabung ke \emph{cluster} lebih dulu sebelum layak
melayani. Masalahnya muncul pada langkah berikut:
\begin{enumerate}[nosep,leftmargin=2em]
  \item Sebuah \emph{pod} baru dinyalakan; ia mengikat \emph{port} dan langsung menandai
        diri ``siap'' (\emph{ready}) semata karena \emph{port}-nya terbuka.
  \item Kubernetes pun mengarahkan pemain ke \emph{pod} itu---padahal \emph{pod} belum
        bergabung ke \emph{cluster}: libcluster (yang menemukan antar-\emph{pod}) belum
        menemukan rekan, dan \emph{registry} Horde belum tersinkron.
  \item Ketika pemain mencari pertandingannya (yang dipegang sebagai \emph{singleton}---satu
        \emph{process} unik untuk pertandingan itu di seluruh \emph{cluster}), \emph{pod} ini
        tak menemukannya (\emph{registry}-nya masih kosong), lalu menjawab ``pertandingan
        tak ada'' atau malah membuat \emph{process} duplikat yang memecah pertandingan.
\end{enumerate}
\noindent Akar masalahnya: ``siap'' ditandai terlalu dini---\emph{port} terbuka belum berarti
siap melayani karena \emph{cluster} belum terbentuk. Diperbaiki M4, yang menahan status
``siap'' sampai \emph{pod} benar-benar tergabung dan \emph{registry} tersinkron.

\section{Churn \emph{Rolling-Deploy}: \emph{Handoff} vs.\ \emph{Grace Period}}
\label{sec:p4}
Pada editor kolaboratif \emph{CoEdit}, tiap dokumen adalah \emph{process} \emph{stateful}
ber-CRDT (\emph{conflict-free replicated data type}---struktur data yang konvergen tanpa
koordinasi) yang di-\emph{handoff} via Horde. Ketika \texttt{SIGTERM} tiba, \emph{handoff}
ratusan \emph{process} memerlukan waktu lebih lama daripada
\texttt{terminationGracePeriodSeconds}; saat \texttt{SIGKILL} datang, dokumen yang belum
berpindah mati bersama suntingan yang belum tersinkron.

\begin{figure}[H]\centering\includegraphics[width=\linewidth,height=0.9\textheight,keepaspectratio]{04-rolling-deploy-sequence}
  \caption{Masalah 4 (sekuens): setelah \texttt{SIGTERM}, \emph{pod} men-\emph{drain} dan
  memulai \emph{handoff} 400 \emph{process}; karena \emph{handoff} ($\approx$45\,s)
  melampaui \emph{grace} (30\,s), \texttt{SIGKILL} memutus sisanya dan state hilang.}\label{fig:p4-seq}\end{figure}
\begin{figure}[H]\centering\includegraphics[width=\linewidth,height=0.9\textheight,keepaspectratio]{04-rolling-deploy-activity}
  \caption{Masalah 4 (aktivitas): perlombaan antara penyelesaian \emph{drain}/\emph{handoff}
  dan tenggat \emph{grace period}; bila tenggat menang, terjadi kehilangan state.}\label{fig:p4-act}\end{figure}
\begin{algorithm}[H]\caption{Masalah 4 --- \emph{handoff} melampaui \emph{grace}.}\label{alg:p4}
\begin{algorithmic}[1]
\Procedure{OnSigterm}{}
  \State \Call{StartGraceTimer}{30}
  \For{$p \in \{400\ \text{proc}\}$} \State \Call{Handoff}{$p$} \Comment{$\approx$45 s} \EndFor
  \If{\textbf{not} \Call{Done}{} \textbf{before} \textit{deadline}} \State \Call{Sigkill}{} \Comment{state hilang} \EndIf
\EndProcedure
\end{algorithmic}\end{algorithm}

\noindent\textbf{Pembahasan.} Diagram sekuens (Gambar~\ref{fig:p4-seq}) menelusuri jalur
terminasi: \emph{Deployment}~$\rightarrow$~kubelet mengirim \texttt{SIGTERM}, \emph{pod}
men-\emph{drain} \emph{endpoint} lalu mengalirkan 400 perintah \emph{handoff} ke node
tersisa; sebuah pesan \texttt{SIGKILL} terjadwal di $t{=}30$\,s memotong aliran yang baru
mencapai $\approx$45\,s. Diagram aktivitas (Gambar~\ref{fig:p4-act}) memodelkan ini sebagai
perlombaan: percabangan ``\emph{handoff} selesai sebelum tenggat?'' menentukan apakah
keluar bersih atau \texttt{SIGKILL}. Algoritma~\ref{alg:p4} memperjelas ketakcocokan
waktu: baris~2 memulai pewaktu 30\,s, perulangan baris~3 melakukan 400 \emph{handoff}
yang totalnya $\approx$45\,s, dan kondisi baris~4 memicu \textsc{Sigkill} begitu tenggat
terlewati---ketiadaan kopling antara durasi \emph{handoff} dan \emph{grace} inilah yang
dijembatani M3.

\subsection{Ilustrasi Sederhana}
\label{subsec:p4-ilustrasi}
Untuk pembaca yang baru mengenal Kubernetes, alur masalah ini ditelusuri langkah demi
langkah dengan angka konkret. Misalkan sebuah \emph{pod} (wadah berisi satu salinan aplikasi
yang dijadwalkan dan dikelola Kubernetes) menjalankan 400 \emph{process} (unit eksekusi
sangat ringan di dalam VM BEAM---mesin virtual tempat Elixir berjalan---sehingga satu
aplikasi lazim menjalankan ratusan ribu sekaligus, berbeda dari proses sistem operasi yang
``berat''), dan tiap \emph{process} memegang \emph{state}: data yang sedang ditahannya di
memori, di sini suntingan dokumen yang belum tersimpan ke basis data. Ketika Kubernetes
memutuskan mematikan \emph{pod} ini (misalnya karena \emph{rolling update}, yakni pembaruan
bertahap yang mengganti \emph{pod} satu per satu ketika ada versi baru aplikasi, perubahan
konfigurasi, atau pembaruan lain), terjadi rangkaian berikut (Tabel~\ref{tab:p4-timeline}):

\begin{enumerate}[nosep,leftmargin=2em]
  \item Pada detik ke-0, Kubernetes mengirim \texttt{SIGTERM} (sinyal yang meminta proses
        berhenti dengan sopan---aplikasi masih diberi kesempatan membereskan diri lebih
        dulu) dan memulai pewaktu \emph{grace} (jatah waktu agar aplikasi sempat selesai
        sebelum dimatikan paksa) selama 30~detik; \emph{pod} lalu berhenti menerima koneksi
        baru (disebut \emph{drain}---menuntaskan permintaan yang sedang berjalan sambil
        menolak yang baru).
  \item \emph{Pod} mulai memindahkan ke-400 \emph{process} beserta \emph{state}-nya ke node
        lain (mesin lain---fisik atau virtual---tempat \emph{pod} dijalankan, satu dari
        banyak node dalam \emph{cluster}) melalui \emph{handoff} (pemindahan \emph{process}
        beserta datanya ke node lain agar tidak ikut hilang saat \emph{pod} mati, dikelola
        oleh Horde---pustaka Elixir yang mengatur \emph{process} secara terdistribusi
        lintas-node dan memindahkannya otomatis ketika sebuah node hilang); karena tiap
        pemindahan harus menyalin \emph{state}, seluruhnya butuh sekitar 45~detik.
  \item Pada detik ke-30 pewaktu \emph{grace} habis, sehingga Kubernetes mengirim
        \texttt{SIGKILL} (sinyal mematikan paksa yang tidak bisa ditolak atau ditunda
        aplikasi) dan \emph{pod} mati seketika---memotong pemindahan yang baru selesai
        $\approx$67\%.
\end{enumerate}

\begin{table}[H]\centering\small
\caption{Garis waktu terminasi pada Bagian~\ref{sec:p4} (\emph{handoff} 45~s vs.\ \emph{grace} 30~s).}
\label{tab:p4-timeline}
\begin{tabular}{@{}p{0.16\linewidth}p{0.78\linewidth}@{}}
\toprule
\textbf{Waktu} & \textbf{Kejadian} \\
\midrule
$t = 0$~s & \texttt{SIGTERM} (minta berhenti sopan); pewaktu \emph{grace} mulai; \emph{pod} berhenti menerima koneksi (\emph{drain}) \\
$t = 0$--$45$~s & \emph{Handoff} 400 \emph{process}+\emph{state} ke node lain ($\approx$9 \emph{process}/detik) \\
$t = 30$~s & Pewaktu \emph{grace} habis $\rightarrow$ \texttt{SIGKILL} (mematikan paksa) $\rightarrow$ \emph{pod} mati seketika \\
Akibat & $\approx$150 \emph{process} yang belum dipindah mati bersama \emph{state}-nya (data hilang) \\
\bottomrule
\end{tabular}
\end{table}

\noindent Akar persoalannya: nilai \emph{grace} (30~detik) ditetapkan \textbf{statis}, tanpa
mengetahui bahwa \emph{handoff} sesungguhnya butuh 45~detik. Perbaikannya---mekanisme M3
(Bab~\ref{ch:mechanisms})---menghitung \emph{grace} secara dinamis dari jumlah \emph{process}
dan laju \emph{handoff}, sehingga tenggat selalu cukup dan tak ada \emph{state} yang hilang.

\section{\emph{Split-Brain}}
\label{sec:p5}
Pada papan peringkat \emph{TopScore}, partisi jaringan memecah \emph{cluster}; tiap sisi
melihat sisi lain sebagai \emph{nodedown} (node distribusi yang terputus) dan mendaftarkan
ulang \emph{singleton}-nya (proses unik se-\emph{cluster} via \texttt{:global}),
sementara semua \emph{probe} lokal Kubernetes tetap lulus. Saat partisi pulih, resolusi
konflik membunuh salah satu duplikat dan membuang akumulasi datanya.

\begin{figure}[H]\centering\includegraphics[width=\linewidth,height=0.9\textheight,keepaspectratio]{05-split-brain-sequence}
  \caption{Masalah 5 (sekuens): partisi memecah \emph{cluster} menjadi dua sisi yang
  sama-sama melayani; \emph{probe} per-\emph{pod} semuanya lulus sehingga Kubernetes tak
  bertindak, lalu resolusi konflik saat pulih membuang state satu duplikat.}\label{fig:p5-seq}\end{figure}
\begin{figure}[H]\centering\includegraphics[width=\linewidth,height=0.9\textheight,keepaspectratio]{05-split-brain-activity}
  \caption{Masalah 5 (aktivitas): siklus \emph{split-brain}---kedua sisi berlanjut dan
  mendaftarkan \emph{singleton}, lalu satu duplikat dibunuh saat pulih (state hilang).}\label{fig:p5-act}\end{figure}
\begin{algorithm}[H]\caption{Masalah 5 --- \emph{split-brain} dan resolusi merusak.}\label{alg:p5}
\begin{algorithmic}[1]
\Procedure{OnPartition}{}
  \State tiap sisi: lihat sisi lain \textit{nodedown}; \Call{Register}{\textit{singleton}}
  \State \textbf{// probe lokal K8s semua lulus $\Rightarrow$ tak bertindak}
\EndProcedure
\Procedure{OnHeal}{}
  \State \Call{Global.ResolveConflict}{} \Comment{bunuh duplikat $\rightarrow$ state hilang}
\EndProcedure
\end{algorithmic}\end{algorithm}

\noindent\textbf{Pembahasan.} Diagram sekuens (Gambar~\ref{fig:p5-seq}) memperlihatkan
dua kelompok \emph{pod} (Partisi~A dan~B) yang, setelah tautan putus, masing-masing
menganggap sisi lain \emph{nodedown} dan terus melayani; di jalur paralel, \emph{probe}
kubelet ke tiap \emph{pod} tetap lulus karena bersifat lokal. Diagram aktivitas
(Gambar~\ref{fig:p5-act}) menyusun siklus hidupnya: dari operasi normal, masuk ke dua
cabang yang sama-sama mendaftarkan \emph{singleton}, lalu menyatu kembali pada peristiwa
\emph{heal}. Algoritma~\ref{alg:p5} memisahkan dua fase: \textsc{OnPartition} (baris~2--3)
menunjukkan kedua sisi mendaftar dan Kubernetes diam; \textsc{OnHeal} (baris~5) memanggil
\textsc{Global.ResolveConflict} yang membunuh satu duplikat sehingga state-nya hilang.
Ketidaksepakatan antara model kesehatan orkestrator (semua hijau) dan model keanggotaan
BEAM (terpecah) inilah yang menjadi sasaran M5.

\subsection{Ilustrasi Sederhana}
\label{subsec:p5-ilustrasi}
Pada \emph{TopScore}:
\begin{enumerate}[nosep,leftmargin=2em]
  \item Jaringan terputus dan memecah \emph{cluster} menjadi dua kelompok \emph{pod} yang
        tak saling melihat; tiap sisi mengira sisi lain mati (\emph{nodedown}).
  \item Tiap sisi tetap melayani dan masing-masing membuat papan-peringkat (\emph{singleton})
        sendiri, sehingga ada dua salinan yang sama-sama menerima skor.
  \item \emph{Probe} Kubernetes ke tiap \emph{pod} tetap lulus---karena \emph{probe} bersifat
        lokal per-\emph{pod}---jadi Kubernetes tak melihat ada yang salah.
  \item Saat jaringan pulih, mekanisme penyelesaian konflik membunuh salah satu duplikat,
        dan skor yang terlanjur dikumpulkannya ikut hilang.
\end{enumerate}
\noindent Akar masalahnya: Kubernetes melihat ``semua sehat'' sementara BEAM melihat
``terpecah''---kedua model tak sepakat. Disasar M5.

\section{Pertanyaan Penelitian}
\label{sec:rq}
\begin{description}[leftmargin=2.6em,style=sameline,nosep]
  \item[RQ1] Komunikasi BEAM$\leftrightarrow$Kubernetes apa yang diperlukan, dan bagaimana
        ia diwujudkan dalam praktik?
  \item[RQ2] Ragam kegagalan apa yang muncul dari interaksi toleransi kesalahan lapis-OTP
        dan lapis-Kubernetes?
  \item[RQ3] Bagaimana dampak toleransi kesalahan berlapis terhadap waktu pemulihan,
        ketersediaan, dan konsistensi state dibandingkan \emph{baseline} satu lapis?
  \item[RQ4] Dapatkah mekanisme koordinasi eksplisit meningkatkan keandalan dibanding
        konfigurasi baku/\emph{ad hoc}?
  \item[RQ5] Kapan pemulihan sebaiknya dimiliki runtime vs.\ orkestrator, dan apakah
        temuannya berlaku umum di luar Elixir?
\end{description}
